\documentclass[conference]{IEEEtran}

\usepackage{graphicx}
\usepackage{amsmath,amssymb}
\usepackage{float}
\usepackage{booktabs}
\usepackage{multirow}

\begin{document}

\title{
    \textbf{Benchmark Report: 502.gcc\_r} \\
    \text{Register Spill Analysis on RISC-V via gem5}
}

\author{
    \IEEEauthorblockN{Lütfullah Burak Kaya}
    \IEEEauthorblockA{
        Ozyegin University \\
        burak.kaya.50711@ozu.edu.tr
    }
    \and
    \IEEEauthorblockN{Ismail Akturk}
    \IEEEauthorblockA{
        Ozyegin University \\
        ismail.akturk@ozyegin.edu.tr
    }
}

\newcommand{\LongDate}{%
    \number\day\space
    \ifcase\month
    \or January\or February\or March\or April\or May\or June%
    \or July\or August\or September\or October\or November\or December%
    \fi
    \space \number\year
}

\twocolumn[
\begin{@twocolumnfalse}
\maketitle
\begin{center}{\small \LongDate}\end{center}
\vspace{1em}
\end{@twocolumnfalse}
]

% ================================================

\begin{abstract}
This report presents the register spill characterization of the
\texttt{502.gcc\_r} benchmark from SPEC CPU2017, executed on a
RISC-V (RV64GC) target using the gem5 TimingSimpleCPU simulator.
A custom SpillDetector was integrated into gem5 to count store-then-load
patterns on the stack within a user-defined Region of Interest (ROI).
The test dataset (\texttt{t1.c}, a minimal C file compiled with
\texttt{-O3 -finline-limit=50000}) was used as the input.
Results show a spill rate of \textbf{6.92\%} within the ROI ---
third highest across all SPEC CPU2017 benchmarks evaluated in this
study --- with a spill-to-load ratio of \textbf{37.5\%}, the highest
observed fraction among all benchmarks. The GCC compilation pipeline
maintains a large set of simultaneously live compiler-internal data
structures on the RISC-V stack, producing severe and persistent
register pressure throughout the \texttt{toplev\_main()} call.
\end{abstract}

% =========================================================
\section{Benchmark Overview}

\texttt{502.gcc\_r} is the SPEC CPU2017 rate variant of GCC 4.5.0
(SPEC-modified), a C compiler that translates C source files into
assembly output. It is one of the largest and most complex integer
workloads in the SPEC suite: the binary encompasses the full GCC
front-end (lexer, parser, semantic analysis), middle-end (SSA
construction, loop optimisations, inlining, register allocation),
and back-end (instruction selection, scheduling, assembly emission).

The \texttt{502.gcc\_r} binary compiled for RISC-V is itself a
RISC-V executable; however, the GCC it embodies was configured to
\emph{target} the x86 architecture. Consequently, the output assembly
files contain x86 instructions, while the register spill measurements
reflect the RISC-V execution of the GCC binary itself.

The RISC-V binary (\texttt{gcc\_r.riscv}) was compiled with gem5
ROI markers inserted around the \texttt{toplev\_main(argc, argv)}
call in \texttt{src/main.c}:

\begin{itemize}
    \item \texttt{m5\_work\_begin(0,0)} --- before \texttt{toplev\_main()}
    \item \texttt{m5\_work\_end(0,0)}   --- after  \texttt{toplev\_main()}
\end{itemize}

\texttt{toplev\_main()} is the entry point to GCC's entire compilation
pipeline and encompasses all parsing, optimisation, and code-generation
passes.

\subsection{Input Datasets}

SPEC CPU2017 provides three dataset sizes for \texttt{gcc\_r}, each
supplying different C source files and compiler option combinations
(Table~\ref{tab:inputs}).

\begin{table}[H]
\centering
\caption{gcc\_r Input Dataset Sizes}
\label{tab:inputs}
\begin{tabular}{lll}
\toprule
\textbf{Dataset} & \textbf{Input file(s)} & \textbf{Compiler options} \\
\midrule
test    & t1.c (69\,B)      & \texttt{-O3 -finline-limit=50000} \\
train   & Various .c files  & Multiple option sets \\
refrate & gcc-pp.c ($>$1\,MB) & Multiple option sets, many runs \\
\bottomrule
\end{tabular}
\end{table}

The \textbf{test} dataset was selected for this study for consistency
with the other evaluated benchmarks. The test input \texttt{t1.c} is a
minimal C file (\texttt{printf("hello world")}) compiled with
aggressive inlining enabled (\texttt{-finline-limit=50000}).
Although the source file is small, the GCC compilation pipeline incurs
substantial fixed overhead per compilation unit: front-end
initialisation, option processing, pass scheduling, and back-end
setup all execute regardless of source size, resulting in 16 million
ROI instructions.

% =========================================================
\section{Simulation Setup}

\subsection{gem5 Configuration}

\begin{itemize}
    \item \textbf{CPU model:}   TimingSimpleCPU
    \item \textbf{ISA:}         RISC-V RV64GC
    \item \textbf{Caches:}      enabled (default L1 I/D + L2)
    \item \textbf{Memory:}      4\,GB simulated physical RAM
    \item \textbf{Spill mode:}  summary only (\texttt{spill\_verbose=False})
\end{itemize}

The full gem5 invocation was:

\begin{verbatim}
./build/RISCV/gem5.opt
  --outdir=benchmarks/gcc/m5out_test
  configs/deprecated/example/se.py
  --cmd=benchmarks/gcc/gcc_r.riscv
  --options="benchmarks/gcc/t1.c
             -O3 -finline-limit=50000
             -o benchmarks/gcc/t1.s"
  --cpu-type=TimingSimpleCPU
  --caches
  --mem-size=4GB
\end{verbatim}

% =========================================================
\section{Simulation Results}

The simulation completed successfully, producing a 369-byte x86
assembly file (\texttt{t1.s}) as expected for a SPEC-configured
x86-targeting GCC binary.

\subsection{Pre-ROI Context}

Before \texttt{toplev\_main()} was entered, the binary executed
GCC's startup phase: dynamic linker resolution, C runtime
initialisation, option parsing, and language front-end registration.
Table~\ref{tab:preroi} shows the global counters at ROI entry.

\begin{table}[H]
\centering
\caption{Global Counters at ROI Entry}
\label{tab:preroi}
\begin{tabular}{lr}
\toprule
\textbf{Metric} & \textbf{Value} \\
\midrule
Total instructions (pre-ROI) & 2,391,942 \\
Total spills detected (pre-ROI) & 79,848 \\
Pre-ROI spill rate & 3.34\% \\
\bottomrule
\end{tabular}
\end{table}

\subsection{ROI Spill Statistics}

Table~\ref{tab:roi} presents the spill statistics collected within
the ROI.

\begin{table}[H]
\centering
\caption{ROI Spill Statistics --- 502.gcc\_r (test, t1.c -O3)}
\label{tab:roi}
\begin{tabular}{lr}
\toprule
\textbf{Metric} & \textbf{Value} \\
\midrule
ROI instructions         & 16,027,045 \\
ROI stores               &  1,955,576 \\
ROI loads                &  2,957,463 \\
\midrule
\textbf{ROI spills}      & \textbf{1,108,838} \\
\textbf{Spill rate}      & \textbf{6.9185\%} \\
\bottomrule
\end{tabular}
\end{table}

\subsection{Timing}

\begin{table}[H]
\centering
\caption{Simulation Timing}
\label{tab:timing}
\begin{tabular}{lr}
\toprule
\textbf{Metric} & \textbf{Value} \\
\midrule
Simulated time           & 0.042\,s \\
Host wall time           & 11.87\,s \\
Total simulated instructions & 20,805,917 \\
\bottomrule
\end{tabular}
\end{table}

% =========================================================
\section{Analysis}

\subsection{Spill Rate Interpretation}

The measured spill rate of \textbf{6.92\%} places \texttt{gcc\_r}
third among all evaluated SPEC CPU2017 benchmarks
(Table~\ref{tab:compare}), behind only \texttt{507.cactuBSSN\_r}
(7.88\%) and \texttt{500.perlbench\_r} (7.28\%).

\begin{table}[H]
\centering
\caption{Cross-Benchmark Spill Rate Comparison}
\label{tab:compare}
\begin{tabular}{lrr}
\toprule
\textbf{Benchmark} & \textbf{ROI Spills} & \textbf{Spill Rate} \\
\midrule
507.cactuBSSN\_r & 3,782,099,438 & 7.8756\% \\
500.perlbench\_r &     1,070,055 & 7.2838\% \\
502.gcc\_r       &     1,108,838 & \textbf{6.9185\%} \\
526.blender\_r   &    61,845,869 & 6.2193\% \\
523.xalancbmk\_r &    18,381,250 & 4.6427\% \\
541.leela\_r     & 1,158,236,271 & 4.4438\% \\
519.lbm\_r       &   274,130,440 & 4.2130\% \\
531.deepsjeng\_r & 1,928,903,656 & 4.1749\% \\
544.nab\_r       &   151,310,616 & 2.2612\% \\
538.imagick\_r   &     2,424,899 & 2.0601\% \\
505.mcf\_r       &   446,113,136 & 1.8263\% \\
508.namd\_r      &   327,718,955 & 1.0269\% \\
\bottomrule
\end{tabular}
\end{table}

\subsection{Store vs.\ Load Balance}

The load count (2.96M) is $1.51\times$ higher than the store count
(1.96M). This load-dominated ratio is distinctive among the evaluated
benchmarks: most floating-point benchmarks exhibit store-heavy or
near-symmetric ratios (2:1 or higher), whereas \texttt{gcc\_r}
is read-heavy. The GCC compilation pipeline reads complex pointer-heavy
data structures --- expression trees, control-flow graphs, SSA use-def
chains, loop nests, and dominance trees --- far more often than it
writes them, as most passes traverse and query structures that were
built in earlier passes.

\subsection{Spill Fraction of Loads}

Within the ROI, the ratio of spills to total loads is:
\[
\frac{1{,}108{,}838}{2{,}957{,}463} \approx 37.5\%
\]
This is \textbf{the highest spill-to-load fraction observed across all
benchmarks in the study}. Approximately 2 in 5 load instructions is a
spill reload. GCC's mid-end optimisation passes simultaneously maintain
live pointers to expression trees, basic-block data, loop information,
alias sets, and per-instruction metadata. With RISC-V providing only
32 general-purpose registers, the compiler of the compiler cannot keep
all these data-structure pointers live across the deep call chains of
GCC's pass manager, resulting in pervasive reloading of previously
spilled pointers.

\subsection{Pre-ROI Spill Density}

The pre-ROI spill rate of 3.34\% ($79{,}848$ spills over
$2{,}391{,}942$ instructions) reflects GCC's startup phase: option
parsing, language front-end registration, and internal table
initialisation. Even before the main compilation pass begins, GCC
exercises complex initialisation code paths with high register pressure.

\subsection{ROI Coverage}

The ROI accounts for $16{,}027{,}045$ out of $20{,}805{,}917$
total simulated instructions:
\[
\frac{16{,}027{,}045}{20{,}805{,}917} \approx 77.0\%
\]
The remaining 23\% is split between the pre-ROI startup phase and
the post-ROI cleanup and output phase.

\subsection{Compiler-Driven Register Pressure}

The extreme spill-to-load ratio of 37.5\% reflects a structural
property of compiler internals: GCC's passes are written to operate
on a rich set of live pointers simultaneously. During a single pass
invocation, the pass manager keeps alive pointers to the current
function's \texttt{cgraph\_node}, its \texttt{function} struct, the
current \texttt{basic\_block}, the current \texttt{gimple} statement,
and various per-pass state structures. Each RISC-V function call in
GCC's highly modular pass infrastructure spills and reloads these
pointers from the stack, as the 32-register RISC-V file cannot hold
them all simultaneously. This structural property is independent of
the input source file size and would persist on larger inputs.

% =========================================================
\section{Conclusion}

The \texttt{502.gcc\_r} benchmark exhibits a register spill rate of
6.92\% within the \texttt{toplev\_main()} ROI on RISC-V RV64GC, the
third highest across all evaluated benchmarks. Most remarkably, its
spill-to-load ratio of 37.5\% is the highest in the study: nearly
2 in 5 load instructions is a spill reload. This distinguishes
\texttt{gcc\_r} from the floating-point workloads, which exhibit
store-dominated memory traffic; \texttt{gcc\_r} is load-heavy,
driven by repeated traversal of complex compiler-internal data
structures. The results demonstrate that GCC's pass infrastructure
generates severe register pressure on RISC-V RV64GC even for trivially
small input files, owing to the large number of simultaneously live
data-structure pointers that exceed the 32-register file capacity.

\end{document}
